\subsection{Classes of Algebras}

We denote by \(\mathrm{Iso}_{\mathrm{K}}\{\mathrm{A},\mathrm{B}\}\) the relation

``A and B are isomorphic K-algebras of finite degree.''

This is an equivalence relation with respect to A and B. Let A be a K-algebra of finite degree; the class of A, denoted by \(\operatorname{cl}(A)\) (or sometimes \(\operatorname{cl}_{\mathrm{K}}(\mathrm{A})\) ), is the class of objects equivalent to A for the relation \(\operatorname{Iso}_{\mathrm{K}}\) (Set Theory, II, §6, No.~9, p.~122). By definition, \(\operatorname{cl}(A)\) is a K-algebra isomorphic to A; two K-algebras of finite degree are in the same class if and only if they are isomorphic.

We denote by \(\mathcal{A}\) the set of pairs \((\mathrm{W},\mu)\) , where \(\mathrm{W}\) is a linear subspace of \(\mathrm{K}^{(\mathbf{N})}\) that is finite-dimensional over \(\mathrm{K}\) and \(\mu\) is a K-bilinear mapping from \(\mathrm{W} \times \mathrm{W}\) to \(\mathrm{W}\) that makes \(\mathrm{W}\) into an (associative and unital) K-algebra. Every K-algebra of finite degree is isomorphic to such an algebra. By loc. cit., the relation

\[
\alpha \text {   is   a   class   of   K - algebras   of   finite   degree }
\]

is therefore collectivizing in \(\alpha\) (Set Theory, II, §1, No.~4, p.~68). We denote the set of classes of K-algebras of finite degree by \(\mathcal{C}_{\mathrm{K}}\) .

PROPOSITION 1. --- The set \(C_{K}\) , endowed with the law of composition given by \((\alpha, \beta) \mapsto \operatorname{cl}(\alpha \otimes_{\mathrm{K}} \beta)\) , is a commutative monoid. The identity element of \(C_{K}\) is the class \(\varepsilon\) of the K-algebra K. Moreover, if A and B are K-algebras of finite degree, then we have the relation

\[
\operatorname{cl} (\mathrm{A} \otimes_ {\mathrm{K}} \mathrm{B}) = \operatorname{cl} (\mathrm{A}) \operatorname{cl} (\mathrm{B}).\tag{1}
\]

Let A, B, and C be K-algebras of finite degree, with respective classes \(\alpha\) , \(\beta\) , and \(\gamma\) . The K-algebras A and \(\alpha\) are isomorphic, as are B and \(\beta\) . Therefore, the K-algebras \(A \otimes_{K} B\) and \(\alpha \otimes_{K} \beta\) are isomorphic, and we have \(\operatorname{cl}(A \otimes_{K} B) = \operatorname{cl}(\alpha \otimes_{K} \beta)\) , which gives formula (1). It follows that \((\alpha\beta)\gamma\) is the class of the K-algebra \((A \otimes_{K} B) \otimes_{K} C\) and \(\alpha(\beta\gamma)\) that of the K-algebra \(A \otimes_{K} (B \otimes_{K} C)\) . Now, these K-algebras are isomorphic (III, §4, No.~1, p.~461), so we have the equality \((\alpha\beta)\gamma = \alpha(\beta\gamma)\) . Analogously, the relation \(\alpha\varepsilon = \varepsilon\alpha = \alpha\) follows from the fact that the K-algebras \(A \otimes_{K} K\) , \(K \otimes_{K} A\) , and A are isomorphic, and the relation \(\alpha\beta = \beta\alpha\) follows from the fact that the algebras \(A \otimes_{K} B\) and \(B \otimes_{K} A\) are isomorphic.

In the set \(C_{K}\) , the relation `` \(\alpha\) and \(\beta\) are Morita equivalent algebras'' is an equivalence relation (VIII, p.~100). It is compatible with the law of composition on \(C_{K}\) by Proposition 13, d) of VIII, p.~111. We denote by \(M_{K}\) the quotient monoid of \(C_{K}\) by this equivalence relation and by \(\varphi\) the canonical homomorphism from \(C_{K}\) to \(M_{K}\) . For every K-algebra A of finite degree, we denote by {[}A{]} the image of \(\operatorname{cl}(A)\) by \(\varphi\) . If A and B are K-algebras of finite degree, then we have {[}A{]} = {[}B{]} if and only if the K-algebras A and B are Morita equivalent; moreover, we have the relation

\[
[ \mathrm{A} \otimes_ {\mathrm{K}} \mathrm{B} ] = [ \mathrm{A} ] [ \mathrm{B} ]\tag{2}
\]

in the monoid \(M_{K}\) .

Lemma 1. --- Let A be a K-algebra of finite degree and D a field of finite degree over K. We have \([A]=[D]\) in \(M_{K}\) if and only if there exists an integer \(n\geqslant1\) such that A is isomorphic to \(\mathbf{M}_{n}(\mathrm{K})\) .

By definition, we have \([A]=[D]\) if and only if there exists an invertible \((A,D)\) -bimodule, that is (VIII, p.~101, Theorem 1), a right vector space V of nonzero finite dimension over D endowed with an isomorphism of K-algebras \(A\to\operatorname{End}_{\mathrm{D}}(\mathrm{V})\) . Lemma 1 follows.

